% GENERATED FILE. Edit canonical lessons, then run npm run generate:books.
% canonical-source-hash: fnv1a64:5853d29edcefb3bb
% canonical-lessons: ML-C163-purikam, ML-C163-kaivella, ML-C163-kanji, ML-C163-payaru, ML-C163-parippu

\chapter{Eyebrow, Palm, Rice gruel, Beans, Dal}
\label{ch:ml-eyebrow-palm-rice-gruel-beans-dal}
\hlchaptermodality{163}

\begin{chapteropening}
\textbf{By the end of this chapter:} \emph{I can say an eyebrow, a palm, rice gruel, beans and dal.}

Complete the last lesson of chapter 163: i can say an eyebrow, a palm, rice gruel, beans and dal.
\end{chapteropening}

\section[purikaṁ]{\ml{പുരികം} (purikaṁ) — an eyebrow}

\label{lesson:ML-C163-purikam}



\begin{quote}
Before the new one: say the Malayalam for a duck, then the Malayalam for an owl.
\end{quote}

\subsection*{You'll want to know: \ml{പുരികം}}
\textbf{\ml{പുരികം}} — \emph{purikaṁ} — \textquotedblleft{}an eyebrow\textquotedblright{}.

\begin{grammarlens}[title={What you've built}]
One more word for this chapter.
\end{grammarlens}

\subsection*{Your turn}
\begin{itemize}
\raggedright
  \item \emph{Say it:} \emph{purikaṁ}
  \item \emph{Say it:} \emph{purikaṁ}, once more
  \item \emph{Say it:} say \emph{purikaṁ}
  \item \emph{Recall:} say the Malayalam for a duck, then the Malayalam for an owl, then say \emph{purikaṁ} again
\end{itemize}

\begin{tcolorbox}[breakable,colback=teal!4,colframe=teal!35!black,title={Before you move on}]
What does \ml{പുരികം} mean? (\textquotedblleft{}An eyebrow\textquotedblright{}.) Say it once more.
\end{tcolorbox}
\section[kaiveḷḷa]{\ml{കൈവെള്ള} (kaiveḷḷa) — a palm (of the hand)}

\label{lesson:ML-C163-kaivella}



\begin{quote}
Before the new one: say the Malayalam for an owl, then the Malayalam for an eyebrow.

\emph{Rātri}, \textquotedblleft{}night\textquotedblright{}: which Hindi word shares its Sanskrit source? (\emph{Rāt}.)
\end{quote}

\subsection*{You'll want to know: \ml{കൈവെള്ള}}
\textbf{\ml{കൈവെള്ള}} — \emph{kaiveḷḷa} — \textquotedblleft{}a palm (of the hand)\textquotedblright{}.

\begin{grammarlens}[title={What you've built}]
One more word for this chapter.
\end{grammarlens}

\subsection*{Your turn}
\begin{itemize}
\raggedright
  \item \emph{Say it:} \emph{kaiveḷḷa}
  \item \emph{Say it:} \emph{kaiveḷḷa}, once more
  \item \emph{Say it:} say \emph{kaiveḷḷa}
  \item \emph{Recall:} say the Malayalam for an owl, then the Malayalam for an eyebrow, then say \emph{kaiveḷḷa} again
\end{itemize}

\begin{tcolorbox}[breakable,colback=teal!4,colframe=teal!35!black,title={Before you move on}]
What does \ml{കൈവെള്ള} mean? (\textquotedblleft{}A palm (of the hand)\textquotedblright{}.) Say it once more.
\end{tcolorbox}
\section[kañji]{\ml{കഞ്ഞി} (kañji) — rice gruel}

\label{lesson:ML-C163-kanji}



\begin{quote}
Before the new one: say the Malayalam for an eyebrow, then the Malayalam for a palm.
\end{quote}

\subsection*{You'll want to know: \ml{കഞ്ഞി}}
\textbf{\ml{കഞ്ഞി}} — \emph{kañji} — \textquotedblleft{}rice gruel\textquotedblright{}.

\begin{grammarlens}[title={What you've built}]
One more word for this chapter.
\end{grammarlens}

\subsection*{Your turn}
\begin{itemize}
\raggedright
  \item \emph{Say it:} \emph{kañji}
  \item \emph{Say it:} \emph{kañji}, once more
  \item \emph{Say it:} say \emph{kañji}
  \item \emph{Recall:} say the Malayalam for an eyebrow, then the Malayalam for a palm, then say \emph{kañji} again
\end{itemize}

\begin{tcolorbox}[breakable,colback=teal!4,colframe=teal!35!black,title={Before you move on}]
What does \ml{കഞ്ഞി} mean? (\textquotedblleft{}Rice gruel\textquotedblright{}.) Say it once more.
\end{tcolorbox}
\section[payaṟŭ]{\ml{പയറ്} (payaṟŭ) — beans}

\label{lesson:ML-C163-payaru}



\begin{quote}
Before the new one: say the Malayalam for a palm, then the Malayalam for rice gruel.

Greet the day in order. (\emph{Suprabhātaṁ}, \emph{śubha madhyāhnaṁ}, \emph{śubha sāyāhnaṁ}, \emph{śubha rātri}.)
\end{quote}

\subsection*{You'll want to know: \ml{പയറ്}}
\textbf{\ml{പയറ്}} — \emph{payaṟŭ} — \textquotedblleft{}beans\textquotedblright{}.

\begin{grammarlens}[title={What you've built}]
One more word for this chapter.
\end{grammarlens}

\subsection*{Your turn}
\begin{itemize}
\raggedright
  \item \emph{Say it:} \emph{payaṟŭ}
  \item \emph{Say it:} \emph{payaṟŭ}, once more
  \item \emph{Say it:} say \emph{payaṟŭ}
  \item \emph{Recall:} say the Malayalam for a palm, then the Malayalam for rice gruel, then say \emph{payaṟŭ} again
\end{itemize}

\begin{tcolorbox}[breakable,colback=teal!4,colframe=teal!35!black,title={Before you move on}]
What does \ml{പയറ്} mean? (\textquotedblleft{}Beans\textquotedblright{}.) Say it once more.
\end{tcolorbox}
\section[parippŭ]{\ml{പരിപ്പ്} (parippŭ) — dal, split lentils}

\label{lesson:ML-C163-parippu}



\begin{quote}
Before the new one: say the Malayalam for rice gruel, then the Malayalam for beans.
\end{quote}

\subsection*{You'll want to know: \ml{പരിപ്പ്}}
\textbf{\ml{പരിപ്പ്}} — \emph{parippŭ} — \textquotedblleft{}dal, split lentils\textquotedblright{}.

\begin{grammarlens}[title={What you've built}]
One more word for this chapter.
\end{grammarlens}

\subsection*{Your turn}
\begin{itemize}
\raggedright
  \item \emph{Say it:} \emph{parippŭ}
  \item \emph{Say it:} \emph{parippŭ}, once more
  \item \emph{Say it:} say \emph{parippŭ}
  \item \emph{Recall:} say the Malayalam for rice gruel, then the Malayalam for beans, then say \emph{parippŭ} again
\end{itemize}

\begin{tcolorbox}[breakable,colback=teal!4,colframe=teal!35!black,title={Before you move on}]
What does \ml{പരിപ്പ്} mean? (\textquotedblleft{}Dal, split lentils\textquotedblright{}.) Say it once more.
\end{tcolorbox}
